\chapter{Estado del Arte}

La salud móvil (mHealth) ha emergido como una línea consolidada de intervención
en el cuidado oncológico, especialmente para la educación y el seguimiento de
pacientes fuera del entorno hospitalario. La evidencia indica que las
plataformas digitales pueden mejorar el conocimiento y la adherencia de los
cuidadores, facilitando el acceso a información estructurada en el momento en
que surge la duda \citep{nair_2017}. En el contexto pediátrico, el cuidador
principal asume un rol observador y registrador de la evolución del paciente,
por lo que disponer de herramientas que ordenen esa información reduce la carga
cognitiva y emocional del seguimiento domiciliario \citep{lewandowska_2022}.

En relación con el registro y seguimiento de síntomas, la literatura destaca la
utilidad de escalas validadas como el Edmonton Symptom Assessment System (ESAS)
y de aplicaciones que permiten el registro longitudinal de signos vitales y
síntomas en el hogar. Estos antecedentes estructurados se constituyen en
insumos valiosos para la consulta clínica, pues permiten al equipo de salud
revisar la evolución del paciente entre controles en lugar de depender
únicamente del relato del cuidador \citep{arnis_2024}. En el mismo sentido, el
uso de sistemas de apoyo a la decisión y triaje en aplicaciones de cuidado
domiciliario pediátrico ha mostrado resultados favorables en la detección
temprana de situaciones que requieren contacto con el equipo de salud
\citep{lau_2025}.

En oncología pediátrica, los programas nacionales y las sociedades científicas
coinciden en la importancia de la educación del cuidador y en la delimitación
estricta de los contenidos: la información entregada debe ser acotada, validada
y no debe sustituir la evaluación profesional \citep{minsal_sf,pohema_seop_sf}.
En Chile, el Programa PINDA y sus centros de referencia definen protocolos de
emergencia y criterios de derivación que constituyen la base clínica sobre la
que la aplicación construye sus reglas de semaforización
\citep{pinda_estadisticas_2026}.

Finalmente, la evidencia respalda que las intervenciones de salud móvil
dirigidas a cuidadores de pacientes pediátricos oncológicos son viables y bien
aceptadas cuando la interfaz es simple, el lenguaje es comprensible y la
información se encuentra disponible incluso sin conexión
\citep{paho_oms_sf,mskcc_sf,idsa_sf,aepap_sf}. Sobre esta base, el proyecto
OnCuidar adopta un enfoque de seguridad por diseño y funcionamiento offline
-first, alineado con las recomendaciones de organismos internacionales de salud
para el uso de tecnologías digitales en pediatría.